\pagebreak
\section{Conclusions}

We have seen how the rings and $k$-reducibility came forth from the five color theorem. D-reducibility expanded upon $k$-reducibility for rings 6 and above. Then C-reducibility expanded upon D-reducibility by avoiding bad colorings. We have worked through these forms of reducibility with the Birkhoff Diamond and Bernhart Diamond as example. For the Bernhart Diamond, we encountered symmetry faults that require further attention.

The ideas of the four color theorem are generalisations of what we did for the five color theorem. Unavoidability of a reducible configuration plays a key role in both proofs. The theory of Kempe-chains came forth from the need to fix ring colorings to become good colorings of a configuration. For this, we had introduced some forms of notation to clarify our future arguments.

The below Figure illustrates how using more advanced forms of reducibility results in less configurations of ring size 6 and above.

\begin{figure}[!h]
    \centering
    \includegraphics[width=1.0\textwidth]{../beamer/images/reduce.eps}
    \caption{Number of reducible configurations on rings $R_6$ and higher in proofs of the four color theorem}.
\end{figure}

It is expected that there should exist more advanced forms of reducibility that essential hone in on the "heart of the four color theorem". Just like how the vertices of $\deg(v)\leq 5$ were at the heart of the five color theorem. 

The computer part of the four color theorem was only needed to evaluate the trivial calculations of finding implying sets of configurations and reducers. The essence of the four color theorem does not lie in these calculations. In the end, reducibility lies at the heart of the four color theorem and the various implementation and calculation issues caused the proof to become so long and infamous.